\subsection{An objective-matched randomized batch challenger}
\label{sec:joint-precision-challenger}
The following calculation addresses precision success, not expected width.
It uses the uncapped exact-total simplex and the original threshold $1/2$;
the bounded-error calculations above define a different information model.
Write
\[
 J_Q(x)=\max_{r=1,\ldots,6}v_{Qr}^{\top}x,
 \qquad x\geq0,\quad\boldsymbol1^\top x=1,
\]
where $v_{Qr}$ contains the normalized coefficient range within each segment
created by batch $Q$.  These ten batches use the same two prefix queries as
the adaptive policy.  A deterministic batch guarantees precision exactly
when every entry of every $v_{Qr}$ is at most $1/2$.

\paragraph{Full-simplex pair-cover certificate.}
For two batches $Q,R$, suppose that for every contrast pair $(r,s)$ a
number $\lambda_{rs}\in[0,1]$ satisfies
\[
 \lambda_{rs}v_{Qr,i}+(1-\lambda_{rs})v_{Rs,i}
 \leq\tfrac12\qquad(i=1,\ldots,6).
\]
Then $J_Q(x)$ and $J_R(x)$ cannot both exceed $1/2$.  Otherwise some $r,s$
would have both linear values strictly above $1/2$, and their indicated
convex combination would be above $1/2$.  Summing the displayed coordinate
inequalities with weights $x_i$ gives the opposite conclusion.  Consequently,
the equal mixture of $Q,R$ succeeds with probability at least $1/2$ on
\emph{every} feasible record.  This includes threshold equality as success.
Each multiplier is obtained exactly by intersecting six rational interval
constraints on $\lambda$; no discretization of records is used.

\paragraph{Full-domain upper bounds from finitely many records.}
We use the six simplex vertices, their 15 pairwise midpoints, and the
uniform record.  Define $A_{zQ}=\boldsymbol1\{J_Q(z)\leq1/2\}$.
For any rational probability distribution $\mu$ on these 22 feasible
records, let $U=\max_Q\sum_z\mu_z A_{zQ}$.  Every randomized batch policy
has some feasible record with success probability at most $U$: its
average success over $\mu$ is at most $U$, so its minimum cannot exceed
that average.  This is the usual finite-game weak-duality argument, not
a claim that these records exhaust the continuum.  The numerical game
solver supplies candidate weights; rational normalization and direct
evaluation of all ten expectations supply the actual checked upper bound.
If the pair-cover guarantee and this upper bound both equal $1/2$, the
full-continuum randomized-batch optimum is therefore exactly $1/2$.
Elsewhere the interval between the lower and upper bound is retained.

\paragraph{Frozen comparison policy.}
Before replaying answers, we first choose a precision-guaranteeing
deterministic batch if one exists.  Otherwise we choose a certified equal
pair if one exists.  Ties minimize worst expected joint width over the six
vertices, then its uniform vertex average, then the batch indices.  If no
pair certificate exists, we retain the saved mean-width mixture and assign
only the trivial precision lower bound zero.  These fallback policies are
not described as precision-optimal.  Selection uses the frozen coefficients
only, with unchanged normalization and budget.

All 36 multipliers for every selected pair and the feasible adverse
weights are exported.  A solver-free rational checker reconstructs the
coefficients from the hashed source rows, verifies all coordinate
inequalities, recomputes every finite-record success indicator, and checks
the lower and upper bounds.  For the complete panels these bounds match
on 468 of 819 original and 212 of 279 external blocks.  The matching value
is $1/2$ on 165 and 57 blocks, respectively; the remaining matches have
value zero or one.  No continuum optimality is inferred for the unmatched
blocks, nor are these exact-total guarantees transferred to the
bounded-error sensitivity models.
